\par
\SourceLeaf{042}
III. Si $Y^{(\mu)}$ est ordinis 0, i. e. numerus, nulla functio indeterminatae $x$
proprie sic dicta ipsas $Y,Y'$ metiri potest: in hoc itaque casu dicendum est, has
functiones divisorem communem non habere.

IV. Excerpamus ex aequationibus nostris penultimam; dein ex hac eliminemus $Y^{(\mu-1)}$ adiumento aequationis antepenultimae; tunc iterum eliminemus
$Y^{(\mu-2)}$ adiumento aequationis praecedentis et sic porro: hoc pacto habebimus
\[
\begin{aligned}
Y^{(\mu)}&=+kY^{(\mu-2)}-k'Y^{(\mu-1)}\\
&=-k'Y^{(\mu-3)}+k''Y^{(\mu-2)}\\
&=+k''Y^{(\mu-4)}-k'''Y^{(\mu-3)}\\
&=-k'''Y^{(\mu-5)}+k''''Y^{(\mu-4)}\\
&\quad\text{etc.}
\end{aligned}
\]
si functiones $k,k',k''$ etc. ex lege sequente formatas supponamus
\[
\begin{aligned}
k&=1\\
k'&=q^{(\mu-2)}\\
k''&=q^{(\mu-3)}k'+k\\
k'''&=q^{(\mu-4)}k''+k'\\
k''''&=q^{(\mu-5)}k'''+k''\\
&\quad\text{etc.}
\end{aligned}
\]
Erit itaque
\[
\pm k^{(\mu-2)}Y\mp k^{(\mu-1)}Y'=Y^{(\mu)}
\]
valentibus signis superioribus pro $\mu$ pari, inferioribus pro impari. In eo itaque
casu, ubi $Y$ et $Y'$ divisorem communem non habent, invenire licet hoc modo
duas functiones $Z,Z'$ indeterminatae $x$ tales, ut habeatur
\[
ZY+Z'Y'=1
\]
V. Haec propositio manifesto etiam inversa valet, puta, si satisfieri potest
aequationi
\[
ZY+Z'Y'=1
\]
ita, ut $Z,Z'$ sint functiones integrae indeterminatae $x$, ipsae $Y$ et $Y'$ certo divisorem communem habere nequeunt.

\SourceLeaf{043}
\SecondArticle{3}
Disquisitio praeliminaris altera circa transformationem functionum symmetricarum versabitur. Sint $a,b,c$ etc. quantitates indeterminatae, ipsarum multitudo $m$, designemusque per $\lambda'$ illarum summam, per $\lambda''$ summam productorum e
binis, per $\lambda'''$ summam productorum e ternis etc., ita ut ex evolutione producti
\[
(x-a)(x-b)(x-c)\, .\, .\, .\, .
\]
oriatur
\[
x^m-\lambda'x^{m-1}+\lambda''x^{m-2}-\lambda'''x^{m-3}+\text{etc.}
\]
Ipsae itaque $\lambda',\lambda'',\lambda'''$ etc. sunt functiones symmetricae indeterminatarum $a,b,c$ etc.,
i. e. tales, in quibus hae indeterminatae eodem modo occurrunt, sive clarius, tales, quae per qualemcunque harum indeterminatarum inter se permutationem non
mutantur. Manifesto generalius, quaelibet functio integra ipsarum $\lambda',\lambda'',\lambda'''$ etc.
(sive has solas indeterminatas implicet, sive adhuc alias ab $a,b,c$ etc. independentes contineat) erit functio symmetrica integra indeterminatarum $a,b,c$ etc.

\SecondArticle{4}
Theorema inversum paullo minus obvium. Sit $\rho$ functio symmetrica indeterminatarum $a,b,c$ etc., quae igitur composita erit e certo numero terminorum
formae
\[
Ma^\alpha b^\beta c^\gamma\, .\, .\, .\, .
\]
denotantibus $\alpha,\beta,\gamma$ etc. integros non negativos, atque $M$ coëfficientem vel determinatum vel saltem ab $a,b,c$ etc. non pendentem (si forte aliae adhuc indeterminatae praeter $a,b,c$ etc. functionem $\rho$ ingrediantur). Ante omnia inter singulos hos terminos ordinem certum stabiliemus, ad quem finem primo ipsas indeterminatas $a,b,c$ etc. ordine certo per se quidem prorsus arbitrario disponemus, e. g. ita, ut $a$ primum locum obtineat, $b$ secundum, $c$ tertium etc. Dein
e duobus terminis
\[
Ma^\alpha b^\beta c^\gamma\, .\, .\, .\, .\quad\text{et}\quad Ma^{\alpha'}b^{\beta'}c^{\gamma'}\, .\, .\, .\, .
\]
priori ordinem altiorem tribuemus quam posteriori, si fit
\[
\begin{gathered}\text{vel}\quad\alpha>\alpha',\quad\text{vel}\quad\alpha=\alpha'\text{ et }\beta>\beta',\\\text{vel}\quad\alpha=\alpha',\ \beta=\beta'\text{ et }\gamma>\gamma',\quad\text{vel etc.}\end{gathered}
\]

\SourceLeaf{044}
i. e. si e differentiis $\alpha-\alpha',\beta-\beta',\gamma-\gamma'$ etc. prima, quae non evanescit, positiva evadit. Quocirca quum termini eiusdem ordinis non differant nisi respectu
coëfficientis $M$, adeoque in terminum unum conflari possint, singulos terminos
functionis $\rho$ ad ordines diversos pertinere supponemus.

Iam observamus, si $Ma^\alpha b^\beta c^\gamma\, .\, .\, .\, .\, .$ sit ex omnibus terminis functionis $\rho$ is,
cui ordo altissimus competat, necessario $\alpha$ esse maiorem, vel saltem non minorem, quam $\beta$. Si enim esset $\beta>\alpha$, terminus $Ma^\beta b^\alpha c^\gamma\, .\, .\, .\, .\, .$, quem functio $\rho$,
utpote symmetrica, quoque involvet, foret ordinis altioris quam $Ma^\alpha b^\beta c^\gamma\, .\, .\, .\, .\, .$
contra hyp. Simili modo $\beta$ erit maior vel saltem non minor quam $\gamma$; porro $\gamma$
non minor quam exponens sequens $\delta$ etc.: proin singulae differentiae $\alpha-\beta$,
$\beta-\gamma,\gamma-\delta$ etc. erunt integri non negativi.

Secundo perpendamus, si e quotcunque functionibus integris indeterminatarum $a,b,c$ etc. productum confletur, huius terminum altissimum necessario
esse ipsum productum e terminis altissimis illorum factorum. Aeque manifestum
est, terminos altissimos functionum $\lambda',\lambda'',\lambda'''$ etc. resp. esse $a,ab,abc$ etc. Hinc
colligitur, terminum altissimum e producto
\[
p=M{\lambda'}^{\alpha-\beta}{\lambda''}^{\beta-\gamma}{\lambda'''}^{\gamma-\delta}\, .\, .\, .\, .\, .
\]
prodeuntem esse $Ma^\alpha b^\beta c^\gamma\, .\, .\, .\, .\, .$; quocirca statuendo $\rho-p=\rho'$, terminus altissimus functionis $\rho'$ certo erit ordinis inferioris quam terminus altissimus functionis $\rho$. Manifesto autem $p$, et proin etiam $\rho'$, fiunt functiones integrae symmetricae ipsarum $a,b,c$ etc. Quamobrem $\rho'$ perinde tractata, ut antea $\rho$, discerpetur in $p'+\rho''$, ita ut $p'$ sit productum e potestatibus ipsarum $\lambda',\lambda'',\lambda'''$ etc. in
coëfficientem vel determinatum vel saltem ab $a,b,c$ etc. non pendentem, $\rho''$ vero
functio integra symmetrica ipsarum $a,b,c$ etc. talis, ut ipsius terminus altissimus pertineat ad ordinem inferiorem, quam terminus altissimus functionis $\rho'$.
Eodem modo continuando, manifesto tandem $\rho$ ad formam $p+p'+p''+p'''$ etc.
redacta, i. e. in functionem integram ipsarum $\lambda',\lambda'',\lambda'''$ etc. transformata erit.

\SecondArticle{5}
Theorema in art. praec. demonstratum etiam sequenti modo enunciare possumus: Proposita functione quacunque indeterminatarum $a,b,c$ etc. integra symmetrica $\rho$, assignari potest functio integra totidem aliarum indeterminatarum
$l',l'',l'''$ etc. talis, quae per substitutiones $l'=\lambda',l''=\lambda'',l'''=\lambda'''$ etc. transeat
\SourceLeaf{045}
in $\rho$. Facile insuper ostenditur, \textit{hoc unico tantum modo fieri posse.} Supponamus
enim, e duabus functionibus diversis indeterminatarum $l',l'',l'''$ etc. puta tum
ex $r$, tum ex $r'$ post substitutiones $l'=\lambda',l''=\lambda'',l'''=\lambda'''$ etc. resultare eandem functionem ipsarum $a,b,c$ etc. Tunc itaque $r-r'$ erit functio ipsarum
$l',l'',l'''$ etc. per se non evanescens, sed quae identice destruitur post illas substitutiones. Hoc vero absurdum esse, facile perspiciemus, si perpendamus, $r-r'$
necessario compositam esse e certo numero partium formae
\[
M{l'}^\alpha{l''}^\beta{l'''}^\gamma\, .\, .\, .\, .\, .
\]
quarum coëfficientes $M$ non evanescant, et quae singulae respectu exponentium
inter se diversae sint, adeoque terminos altissimos e singulis istis partibus prodeuntes exhiberi per
\[
Ma^{\alpha+\beta+\gamma+\text{etc.}}b^{\beta+\gamma+\text{etc.}}c^{\gamma+\text{etc.}}\, .\, .\, .\, .\, .
\]
et proin ad ordines diversos referendos esse, ita ut terminus absolute altissimus
nullo modo destrui possit.

Ceterum ipse calculus pro huiusmodi transformationibus pluribus compendiis insigniter abbreviari posset, quibus tamen hoc loco non immoramur, quum
ad propositum nostrum sola transformationis possibilitas iam sufficiat.

\SecondArticle{6}
Consideremus productum ex $m(m-1)$ factoribus
\[
\begin{aligned}
&(a-b)(a-c)(a-d)\, .\, .\, .\, .\, .\\
\times{}&(b-a)(b-c)(b-d)\, .\, .\, .\, .\, .\\
\times{}&(c-a)(c-b)(c-d)\, .\, .\, .\, .\, .\\
\times{}&(d-a)(d-b)(d-c)\, .\, .\, .\, .\, .\\
&\text{etc.}
\end{aligned}
\]
quod per $\pi$ denotabimus, et, quum indeterminatas $a,b,c$ etc. symmetrice involvat, in formam functionis ipsarum $\lambda',\lambda'',\lambda'''$ etc. redactum supponemus. Transeat haec functio in $p$, si loco ipsarum $\lambda',\lambda'',\lambda'''$ etc. resp. substituuntur $l',l'',l'''$ etc.
His ita factis, ipsam $p$ vocabimus \textit{determinantem} functionis
\[
y=x^m-l'x^{m-1}+l''x^{m-2}-l'''x^{m-3}+\text{etc.}
\]
